% Generated draft: requires mathematical and source review.
\documentclass[11pt]{article}
\usepackage[margin=25mm]{geometry}
\usepackage{fontspec}
\setmainfont{Latin Modern Roman}
\usepackage{amsmath,amssymb,mathtools,unicode-math}
\setmathfont{Latin Modern Math}
\usepackage{enumitem,needspace,xurl,hyperref}
\setlength{\parindent}{0pt}
\setlength{\parskip}{5pt}
\setlength{\emergencystretch}{3em}
\newcommand{\N}{\mathbb{N}}
\newcommand{\Z}{\mathbb{Z}}
\newcommand{\Q}{\mathbb{Q}}
\newcommand{\R}{\mathbb{R}}
\newcommand{\C}{\mathbb{C}}
\newcommand{\F}{\mathbb{F}}
\newcommand{\A}{\mathbb{A}}
\newcommand{\PP}{\mathbb{P}}
\newcommand{\E}{\mathbb{E}}
\newcommand{\eps}{\varepsilon}
\newcommand{\dd}{\,\mathrm{d}}
\begin{document}
\section*{UnsolvedMath ID: 30005493}
\textbf{Source problem code:} catalogue entry
\subsection{Definitions and mathematical statement}
A finitely presented group is a group admitting a presentation $\langle g_1,\ldots,g_r\mid R_1,\ldots,R_s\rangle$ with finitely many generators and relations. A complex algebraic singularity $(X,x)$ is understood here as the germ of a complex algebraic variety at a singular point $x$; its dimension is its local complex dimension.

For each dimension $n$ under consideration, let $\mathcal C_n$ denote a specified class of $n$-dimensional complex algebraic singularities. The supplied formulation calls these classes ``natural'' but does not specify their members or the range of classes intended. It also does not define the notation $\pi_1(X;x)$.

One possible topological convention is the fundamental group of a sufficiently small punctured analytic neighborhood. Under this convention, choose a local embedding of a representative of $(X,x)$ into $\C^N$, let $B_\varepsilon(x)$ be a sufficiently small Euclidean ball, and choose a base point $b\in (X\cap B_\varepsilon(x))\setminus\{x\}$. The corresponding group is
\[
\pi_1\bigl((X\cap B_\varepsilon(x))\setminus\{x\},b\bigr).
\]
This is a conditional interpretation, not a convention established by the supplied record. If the punctured neighborhood is disconnected, a component and base point must be specified. Using the regular locus of a neighborhood, rather than deleting only $x$, is another interpretation that must be distinguished when the singularity is not isolated.

Once the classes and the meaning of $\pi_1(X;x)$ have been fixed, the requested classification is the determination, up to group isomorphism, of
\[
\mathscr G(\mathcal C_n)
:=\bigl\{[G]\;\bigm|\;G\text{ is finitely presented and there exists }(X,x)\in\mathcal C_n
\text{ with }G\cong\pi_1(X;x)\bigr\}
\]
for every such $n$ and class $\mathcal C_n$. Equivalently, give necessary and sufficient conditions on a finitely presented group $G$ for this realization to exist.

No normality, isolatedness, irreducibility, or particular singularity class is explicitly imposed by the supplied statement. The original wording, ``What fundamental groups can we find among the singularities in Cn?'', does not by itself resolve whether ``Cn'' denotes $\mathcal C_n$ or an ambient affine space $\C^n$. Thus the displayed target faithfully expresses the supplied clean formulation, but a fully specified classification problem requires resolution of these qualifications.
\subsection{Short English statement}
For each specified class of complex algebraic singularities of a given dimension, characterize the finitely presented groups realized as their local fundamental groups. The intended classes and definition of local fundamental group remain unspecified; the original notation ``Cn'' also needs clarification.
\subsection{Sources}
\begin{itemize}
\item Local UnsolvedMath registry, permanent ID $30005493$, source problem code \texttt{catalogue entry}, title: \emph{Local Fundamental Groups of Algebraic Singularities}. The supplied record belongs to the Oberwolfach Reports Open Problems collection and gives the proposed year as 2023; no individual proposer is identified.
\item \emph{Komplexe Analysis -- Differential and Algebraic methods in Kähler spaces} (2023), Oberwolfach Reports, pp.~965--1030, DOI: \url{https://doi.org/10.4171/owr/2023/18}. The supplied original-statement excerpt identifies p.~1007. The record describes Oberwolfach Reports as documenting mathematical workshops and their open-problem discussions.
\item Supplied literature reference, labeled ``Kapovich and Kollar'': \url{https://doi.org/10.1090/S0894-0347-2014-00807-9}.
\item Supplied literature reference, labeled ``Low-dimensional lc singularities'': \url{https://arxiv.org/abs/2302.11790}.
\item Supplied literature reference, labeled ``OWR 2023 report'': \url{https://ems.press/content/serial-article-files/47013}.
\item Literature assessment dated exactly \texttt{2026-08-22}: the supplied curation labels the problem ``partially solved'' and states that normal singularities and several special classes have strong fundamental-group realizability theorems, while the broad classification across singularity classes remains incomplete. This is supplied dated curation derived from the OWR record, not an independently verified assessment; time-sensitive status and cited claims require independent verification.
\item The registry labels the statement \texttt{corrected\_verified} and reports that it was checked against the official Oberwolfach report, preserving its hypotheses, notation, and open target. This verification claim is supplied provenance and has not been independently checked here.
\end{itemize}
\end{document}
